\section{Card models}

Card models are useful because the phrase ``draw cards'' is incomplete. It does
not say whether we record one card or several cards, whether order matters,
whether cards are replaced, or whether the final object is a hand, a sequence,
or only a summary such as the number of aces. Each of these decisions changes
the sample space.

In this section let \(D\) denote a standard deck of \(52\) distinct cards. The
cards are distinct because a card is determined by both rank and suit. For
example, the ace of hearts and the ace of spades are different cards.

\subsection*{Drawing one card}

Suppose one card is drawn from a well-shuffled deck and the exact card is
recorded. The sample space is
\[
\Omega=D.
\]
It has
\[
|\Omega|=52
\]
elementary outcomes. If the deck is well shuffled and no card is favoured, the
uniform model assigns
\[
P(c)=\frac1{52}
\]
for every card \(c\in D\).

Events are subsets of the deck. For example, the event that the card is an ace
is
\[
A=\{\text{ace of hearts},\text{ace of diamonds},\text{ace of clubs},
\text{ace of spades}\}.
\]
Therefore
\[
P(A)=\frac4{52}=\frac1{13}.
\]

The event that the card is a heart contains \(13\) cards, so
\[
P(\text{heart})=\frac{13}{52}=\frac14.
\]
The event that the card is a face card contains the jacks, queens, and kings in
four suits, hence \(12\) cards, so
\[
P(\text{face card})=\frac{12}{52}=\frac3{13}.
\]

\begin{remarkbox}
The sample space contains individual cards, not ranks. If we record the exact
card, then there are \(52\) elementary outcomes. The event ``ace'' is not one
elementary outcome; it is a set of four elementary outcomes.
\end{remarkbox}

\subsection*{A five-card hand as a set}

Now suppose five cards are dealt and only the final hand is recorded. The order
in which the cards arrived is ignored. Then an elementary outcome is a five-card
subset of the deck:
\[
S\subseteq D,\qquad |S|=5.
\]
The sample space is
\[
\Omega=\{S\subseteq D: |S|=5\}.
\]
Its size is
\[
|\Omega|=\binom{52}{5}.
\]

This count appears because a hand is a set. We choose which five cards belong to
the hand, and the order of listing those cards does not create a new hand. Thus
\[
\{A\heartsuit,K\clubsuit,7\diamondsuit,3\spadesuit,2\heartsuit\}
\]
is the same hand no matter in which order we write its elements.

If the deck is well shuffled and every five-card hand is equally likely, then
the uniform formula applies:
\[
P(E)=\frac{|E|}{\binom{52}{5}}
\]
for any event \(E\) consisting of five-card hands.

\begin{examplebox}
Let \(E\) be the event that the hand contains exactly one ace. To construct such
a hand, choose one of the four aces and choose the remaining four cards from the
\(48\) non-aces. Hence
\[
|E|=\binom41\binom{48}{4}.
\]
Therefore
\[
P(E)=\frac{\binom41\binom{48}{4}}{\binom{52}{5}}.
\]
This is a probability of an event, not of a single hand. The event contains many
hands.
\end{examplebox}

\subsection*{Five cards drawn in order}

Suppose instead that five cards are drawn one after another and the order is
recorded. Then an elementary outcome is a sequence
\[
(c_1,c_2,c_3,c_4,c_5),
\]
where all five cards are distinct. The sample space is
\[
\Omega_{\mathrm{seq}}=\{(c_1,c_2,c_3,c_4,c_5): c_i\in D,\ c_i\neq c_j
\text{ for }i\neq j\}.
\]
Its size is
\[
52\cdot 51\cdot 50\cdot 49\cdot 48.
\]
The first position has \(52\) choices, the second has \(51\) choices because one
card has already been used, and so on.

This sample space is different from the five-card hand space. A hand ignores
order; a sequence records order. The same five cards can appear in
\[
5!
\]
different orders. Therefore
\[
52\cdot 51\cdot 50\cdot 49\cdot 48=5!\binom{52}{5}.
\]

\begin{remarkbox}
The phrase ``draw five cards'' is not enough to determine the sample space. If
order is ignored, the elementary outcome is a set. If order is recorded, the
elementary outcome is a sequence. Both models are legitimate, but they answer
different questions.
\end{remarkbox}

\subsection*{The same event in two sample spaces}

Consider the statement
\[
\text{the five cards contain exactly one ace.}
\]
In the hand model, the event is a set of five-card subsets. In the sequence
model, the event is a set of ordered five-tuples. The verbal statement is the
same, but the formal event lives in a different sample space.

In the hand model we already counted
\[
\binom41\binom{48}{4}
\]
favourable hands.

In the ordered model, we may count favourable sequences directly. First choose
which position contains the ace: \(5\) choices. Choose the ace: \(4\) choices.
Then fill the remaining four ordered positions with non-aces without repetition:
\[
48\cdot47\cdot46\cdot45
\]
choices. Thus the number of favourable sequences is
\[
5\cdot4\cdot48\cdot47\cdot46\cdot45.
\]
The total number of ordered sequences is
\[
52\cdot51\cdot50\cdot49\cdot48.
\]
So the probability in the ordered model is
\[
\frac{5\cdot4\cdot48\cdot47\cdot46\cdot45}
{52\cdot51\cdot50\cdot49\cdot48}.
\]
This value is equal to
\[
\frac{\binom41\binom{48}{4}}{\binom{52}{5}},
\]
because both models describe the same random dealing mechanism at different
levels of detail.

The equality is not magic. Each unordered hand corresponds to exactly \(5!\)
ordered sequences. The factor \(5!\) appears in both the favourable count and
the total count, so it cancels.

\subsection*{Drawing with replacement}

A different model appears if each drawn card is returned to the deck before the
next draw. Suppose five draws are made, the exact card is recorded each time,
and replacement occurs after every draw. Then the same card may appear several
times in the recorded sequence. The sample space is
\[
\Omega_{\mathrm{rep}}=D^5.
\]
Its size is
\[
|\Omega_{\mathrm{rep}}|=52^5.
\]
If the deck is well shuffled before each draw, every sequence in \(D^5\) has
probability
\[
\frac1{52^5}.
\]

This is not the same as drawing five cards without replacement. For example,
under replacement the sequence
\[
(A\heartsuit,A\heartsuit,7\clubsuit,2\diamondsuit,K\spadesuit)
\]
is possible. Without replacement it is impossible, because the same physical
card cannot be drawn twice.

\begin{examplebox}
In five draws with replacement, what is the probability that all five cards are
aces?

Each draw has probability \(4/52=1/13\) of producing an ace. Since replacement
makes the same probability available at each draw, the probability is
\[
\left(\frac{4}{52}\right)^5=\left(\frac1{13}\right)^5.
\]
Equivalently, there are \(4^5\) ace-only sequences and \(52^5\) total sequences,
so
\[
P(\text{all aces})=\frac{4^5}{52^5}.
\]
\end{examplebox}

Without replacement, the probability that five drawn cards are all aces is zero,
because the deck contains only four aces. The change in the physical procedure
changes the sample space and therefore changes the probabilities.

\subsection*{At least one ace}

Card models are also useful for showing the complement method. In a five-card
hand, let
\[
A=\{\text{the hand contains at least one ace}\}.
\]
Counting hands with at least one ace directly would require cases: exactly one
ace, exactly two aces, exactly three aces, exactly four aces. It is often simpler
to count the complement.

The complement is
\[
A^c=\{\text{the hand contains no aces}\}.
\]
There are \(48\) non-aces, so
\[
|A^c|=\binom{48}{5}.
\]
Therefore
\[
P(A^c)=\frac{\binom{48}{5}}{\binom{52}{5}},
\]
and
\[
P(A)=1-P(A^c)=1-\frac{\binom{48}{5}}{\binom{52}{5}}.
\]

The complement method is not a trick. It is event algebra: instead of counting
all hands that satisfy the event, we count the hands that fail it and subtract
from the whole space.

\subsection*{What card models teach}

Card models emphasize that the same everyday phrase can hide several different
mathematical structures.

If one card is drawn and the exact card is recorded, then \(\Omega=D\). If five
cards are recorded as a hand, then \(\Omega\) is a set of five-element subsets.
If five cards are recorded in order without replacement, then \(\Omega\) is a
set of sequences with no repeated card. If cards are drawn with replacement,
then \(\Omega=D^5\), and repeated cards become possible.

\begin{remarkbox}
Before solving a card problem, always identify whether the elementary outcome is
a single card, a set of cards, an ordered sequence of cards, or a summary of the
cards. Most card-counting errors are errors about the sample space, not errors
about arithmetic.
\end{remarkbox}
